\documentclass[conference]{IEEEtran}
\IEEEoverridecommandlockouts

% Codificacao e idioma (Portugues Brasileiro)
\usepackage[utf8]{inputenc}
\usepackage[T1]{fontenc}
\usepackage[brazil]{babel}

% Pacotes padrao IEEE e matematicos
\usepackage{cite}
\usepackage{amsmath,amssymb,amsfonts}
\usepackage{algorithmic}
\usepackage{graphicx}
\usepackage{textcomp}
\usepackage{xcolor}
\usepackage{booktabs}
\usepackage{microtype}
\usepackage{url}
\usepackage{tikz}
\usetikzlibrary{arrows.meta,positioning}

% Diretorio padrao de figuras
\graphicspath{{figuras/}}

% Caso queira os rotulos em portugues, descomente as linhas abaixo:
% \renewcommand{\abstractname}{Resumo}
% \renewcommand{\IEEEkeywordsname}{Palavras-chave}

\begin{document}

\title{Comparação entre Post-Training Quantization e Quantization-Aware Training de 8 a 2 Bits no Reconhecimento de Palavras-Chave}

\author{\IEEEauthorblockN{Pedro Hugo Passos da Silva Carlos}
\IEEEauthorblockA{\textit{Departamento de Engenharia Elétrica e de Computação} \\
\textit{Programa de Pós-Graduação em Engenharia Elétrica e Computação (PPGEEC)} \\
\textit{Universidade Federal da Bahia (UFBA)} \\
Salvador, Bahia, Brasil \\
pedrohugo.psc@gmail.com}
% Caso deseje adicionar coautor(es) ou orientador, basta usar \and seguido de outro \IEEEauthorblockN{...}
}

\maketitle

\begin{abstract}
% TODO: Redigir o resumo final apos concluir os resultados e conclusoes (nao usar simbolos, notas de rodape ou formulas matematicas no abstract).
Este artigo apresenta uma comparação sistemática entre as técnicas \textit{Post-Training Quantization} (PTQ) e \textit{Quantization-Aware Training} (QAT), incluindo precisões \textit{sub-8-bit} (4 e 2 bits), aplicadas ao reconhecimento de palavras-chave (\textit{keyword spotting}). Tomando como referência o \textit{pipeline} embarcado no ESP32-S3 via Edge Impulse, o modelo 1D-CNN com entrada MFCC é reproduzido em PyTorch sob quantização simulada para isolar o efeito das estratégias PTQ (\textit{Min-Max} e \textit{MSE}) e QAT (Padrão, LSQ, LSQ+ e LSQ+ Híbrido).
\end{abstract}

\begin{IEEEkeywords}
Keyword Spotting, TinyML, Quantização, PTQ, QAT, LSQ, Sub-8-Bit.
\end{IEEEkeywords}

% Secoes do artigo (edite cada arquivo dentro da pasta secoes/)
\section{Introdução}

Em sistemas embarcados, o uso da tecnologia de reconhecimento de palavras-chave (\textit{keyword spotting}) é amplamente utilizado nos dias atuais, permitindo que o usuário acione dispositivos sem o uso das mãos. As expressões mais populares são ``OK Google'', ``Alexa'' e ``Hey Siri''. Assim, esses dispositivos ficam em escuta contínua até que a palavra-chave seja identificada, para que possam se conectar à internet e executar alguma ação solicitada pelo usuário~\cite{carlos2025}.

Para que esses dispositivos possam reconhecer essas palavras, é necessário embarcar algoritmos de Aprendizagem de Máquina (AM). Todavia, há o desafio de inserir esses modelos nesses \textit{hardwares} (microcontroladores), devido aos seus recursos limitados de memória e processamento. Por conta disso, surge o \textit{Tiny Machine Learning} (TinyML), focado em reduzir e comprimir os modelos de AM~\cite{carlos2025,pereira2023}.

Existem várias técnicas a fim de reduzir os modelos de AM, entre elas, destaca-se a quantização, que mapeia os pesos, vieses e ativações de ponto flutuante (FP32) para representações inteiras de baixa precisão~\cite{ravaglia2021,zhao2021}. Para quantizar um modelo, há duas abordagens principais: o \textit{Post-Training Quantization} (PTQ) e o \textit{Quantization-Aware Training} (QAT). O primeiro quantiza os parâmetros após o treinamento do modelo, utilizando estratégias de calibração como \textit{Min-Max} ou \textit{Mean Squared Error} (MSE)~\cite{zhao2021}. Já no segundo, durante o treinamento, os parâmetros são quantizados de forma simulada (\textit{fake quantization}), permitindo que a rede se ajuste por meio do \textit{Straight-Through Estimator} (STE) considerando o erro de discretização e a acurácia~\cite{bhalgat2020,lin2023}. Embora a quantização para 8~bits (compressão de $4\times$ em relação ao FP32) já apresente resultados satisfatórios com PTQ~\cite{carlos2025,ravaglia2021}, a redução para precisões \textit{sub-8-bit} --- como 4~bits ($8\times$) e 2~bits ($16\times$) --- tende a degradar severamente o desempenho sem reajuste dos parâmetros~\cite{bhalgat2020,lin2023}. Nesse cenário, variantes do QAT baseadas em gradiente, como o \textit{Learned Step Size Quantization} (LSQ) e o LSQ+, formulam a quantização como um problema de otimização conjunta, aprendendo também o passo de quantização e o deslocamento (\textit{offset}) durante o treinamento, diferentemente de abordagens que mantêm esses parâmetros fixos~\cite{bhalgat2020}.

Assim, o presente artigo tem como objetivo investigar essas duas abordagens de quantização tomando como base o trabalho ``Sistema Embarcado para Reconhecimento de Palavras usando ESP32-S3 e TinyML'', de Carlos~\cite{carlos2025}. Nesse estudo, o autor utilizou a plataforma Edge Impulse desde o pré-processamento dos dados até a implantação no microcontrolador, aplicando apenas a técnica PTQ para converter os parâmetros de ponto flutuante para inteiros de 8~bits. Diante disso, em vez de utilizar o Edge Impulse, este trabalho reproduz em PyTorch o mesmo \textit{pipeline} de extração de características MFCC e a arquitetura 1D-CNN (1.974 parâmetros), mantendo idênticas as parametrizações originais para focar exclusivamente na etapa de quantização. Com isso, busca-se investigar se há vantagem em utilizar o QAT em sua forma padrão e em suas variantes LSQ, LSQ+ e LSQ+ Híbrido frente ao PTQ (\textit{Min-Max} e \textit{MSE}). Para tanto, pesos e ativações são quantizados para 8~bits, 4~bits e 2~bits, comparando-se os resultados por meio das métricas de acurácia, função de perda (\textit{loss}), AUC e \textit{F1-score} ponderado.

O restante deste artigo está organizado da seguinte forma: a Seção~II apresenta a fundamentação teórica sobre quantização e formulação de otimização das técnicas PTQ e QAT avaliadas; a Seção~III descreve a metodologia, o conjunto de dados e a configuração experimental em PyTorch; a Seção~IV analisa os resultados comparativos nas precisões de 8, 4 e 2~bits; por fim, a Seção~V apresenta as conclusões e trabalhos futuros.

\section{Fundamentação Teórica}

Em redes neurais, os pesos e as ativações são convencionalmente treinados em ponto flutuante de 32~bits. Assim, ao converter esses parâmetros de 32~bits para representações inteiras
de baixa precisão (8, 4 e 2~bits), tanto o uso de memória quanto o custo computacional das multiplicações matriciais diminuem consideravelmente, permitindo que o \textit{hardware}   
execute essas operações com maior eficiência~\cite{nagel2021}.

\subsection{Quantização Linear Uniforme}

Nesse contexto, os métodos de quantização podem ser categorizados em não-lineares e lineares. Embora a abordagem não-linear apresente bons resultados em baixas larguras de bits, ela
exige conversões adicionais para mapear os valores quantizados, como o uso de uma tabela de consulta (\textit{look-up table}). Por isso, adota-se geralmente a quantização linear, por
dispensar essas etapas extras e ser mais adequada à implementação em \textit{hardware}. Essa abordagem subdivide-se em quantização simétrica e assimétrica~\cite{lin2023}.

A quantização simétrica é definida a seguir~\cite{bhalgat2020,nagel2021}:
\begin{align}
\bar{x} &= \operatorname{clamp}\!\left(\left\lfloor \frac{x}{s} \right\rceil,\, n,\, p\right), \label{eq:quant_sim_q} \\
\hat{x} &= \bar{x} \cdot s, \label{eq:quant_sim_dq}
\end{align}
em que $x$ representa o valor real de entrada (peso ou ativação), $s > 0$ é o fator de escala (passo de quantização), $\lfloor \cdot \rceil$ denota o arredondamento para o inteiro mais próximo, $\bar{x}$ é a representação inteira quantizada e $\hat{x}$ é o valor reconstruído (desquantizado). A função $\operatorname{clamp}(\cdot,\, n,\, p)$ satura os níveis no intervalo inteiro de $b$~bits: $[0,\, 2^{b}-1]$ para valores sem sinal (\textit{unsigned}, como ativações pós-ReLU) e $[-2^{b-1},\, 2^{b-1}-1]$ para valores com sinal (\textit{signed}, como pesos)~\cite{bhalgat2020,nagel2021}.

Quando pesos ou ativações apresentam distribuições assimétricas ou não centradas em zero, adota-se a quantização assimétrica (ou uniforme afim), que introduz um parâmetro de deslocamento (\textit{offset}) $\beta$ para ajustar melhor a grade discreta à faixa real de valores~\cite{bhalgat2020,nagel2021}:
\begin{align}
\bar{x} &= \operatorname{clamp}\!\left(\left\lfloor \frac{x - \beta}{s} \right\rceil,\, n,\, p\right), \label{eq:quant_assim_q} \\
\hat{x} &= \bar{x} \cdot s + \beta. \label{eq:quant_assim_dq}
\end{align}
Nessas equações, $\beta$ desloca o intervalo de quantização para cobrir apenas a região onde os dados realmente variam, evitando desperdiçar níveis discretos~\cite{bhalgat2020} --- formulação equivalente à do ponto zero inteiro (\textit{zero-point}) $z$ adotada por Nagel et al.~\cite{nagel2021}, com $\beta = -s \cdot z$. Assim, a quantização assimétrica aproveita melhor os $2^b$ níveis disponíveis em distribuições deslocadas do zero, mas exige operações extras de soma e subtração na inferência para compensar $\beta$, ao contrário do caso simétrico ($\beta = 0$)~\cite{lin2023,nagel2021}.

Em ambas as abordagens, os limites inteiros $[n,\, p]$ correspondem no domínio real aos limiares de corte $[q_{\min},\, q_{\max}]$, dados por $q_{\min} = n\cdot s + \beta$ e $q_{\max} = p\cdot s + \beta$. A escolha de $s$ envolve um equilíbrio entre dois tipos de erro: valores que ultrapassam $[q_{\min},\, q_{\max}]$ são cortados (erro de \textit{clipping}), enquanto os valores internos sofrem erro de arredondamento (de no máximo $s/2$). Aumentar $s$ amplia o alcance $[q_{\min},\, q_{\max}]$ e reduz os cortes, porém aumenta o espaçamento entre os níveis e, consequentemente, o erro de arredondamento~\cite{nagel2021}.

\subsection{Post-Training Quantization (PTQ): Min-Max e MSE}

Os algoritmos de \textit{Post-Training Quantization} (PTQ) partem de uma rede neural já treinada em ponto flutuante (FP32) e a convertem diretamente para ponto fixo,
sem exigir o retreinamento do modelo. Por dispensarem o ajuste complexo de hiperparâmetros e operarem com poucos ou nenhum dado adicional, os métodos de PTQ são computacionalmente eficientes e fáceis de aplicar. Contudo, para preservar o desempenho do modelo, é necessário definir     
adequadamente os intervalos de quantização de cada ponto da rede~\cite{nagel2021}.

A definição do intervalo de quantização consiste em estabelecer os limiares reais de corte, $q_{\min}$ e $q_{\max}$, que determinam o fator de escala $s$:
\begin{equation}
s = \frac{q_{\max} - q_{\min}}{p - n},
\label{eq:escala_s}
\end{equation}
Essa escolha busca equilibrar o erro de saturação (\textit{clipping}) e o erro de arredondamento, idealmente visando minimizar o impacto sobre a função de perda da tarefa. Entretanto, 
como o PTQ evita o treinamento ponta a ponta (\textit{end-to-end}), os métodos de calibração substituem a perda global da tarefa pela otimização de funções de custo locais para cada quantizador,      
garantindo rapidez computacional~\cite{nagel2021}. Entre os critérios mais comuns destacam-se o \textit{Min-Max} e a minimização do \textit{Mean Squared Error} (MSE), descritos a     
seguir.

O critério \textit{Min-Max} tem como objetivo cobrir toda a faixa dinâmica dos valores (pesos ou ativações). Assim, os limiares reais de quantização são definidos da seguinte         
maneira~\cite{nagel2021}:
\begin{align}
q_{\min} &= \min V, \label{eq:ptq_minmax_min} \\
q_{\max} &= \max V, \label{eq:ptq_minmax_max}
\end{align}
em que $V$ denota o conjunto de valores (pesos ou ativações) a ser quantizado. Essa abordagem elimina o erro de corte (\textit{clipping}). Todavia, torna-se sensível a valores        
extremos (\textit{outliers}), pois um único valor distante amplia o intervalo $[q_{\min},\, q_{\max}]$, aumentando o passo $s$ e causando erro excessivo de arredondamento nos valores 
comuns~\cite{nagel2021}.

Já o critério MSE busca aliviar o impacto dos \textit{outliers} para atenuar o erro de arredondamento~\cite{nagel2021}. Para isso, determinam-se os limiares $q_{\min}$ e $q_{\max}$
que minimizam o erro quadrático médio entre os conjuntos de valores (pesos ou ativações) originais e quantizados:                                                                      
\begin{equation}
\arg\min_{q_{\min},\, q_{\max}} \left\| V - \hat{V}(q_{\min},\, q_{\max}) \right\|_F^2,
\label{eq:ptq_mse}
\end{equation}    
em que $\hat{V}(q_{\min},\, q_{\max})$ denota a versão quantizada de $V$ e $\|\cdot\|_F$ representa a norma de Frobenius. Desse modo, busca-se um intervalo que represente melhor a
maioria dos valores, equilibrando o erro de saturação (\textit{clipping}) e o erro de arredondamento~\cite{nagel2021}. Esse problema de otimização é comumente resolvido por métodos de
busca numérica, como a busca em grade (\textit{grid search})~\cite{nagel2021}.

\subsection{Quantization-Aware Training (QAT): STE, LSQ e LSQ+}

A técnica de PTQ, por não envolver o retreinamento da rede, apresenta bom desempenho em quantizações de 8~bits; todavia, ao reduzir a precisão para 4 ou 2~bits, sua acurácia tende a
cair acentuadamente~\cite{bhalgat2020,nagel2021}. Para esses casos, emprega-se o \textit{Quantization-Aware Training} (QAT), no qual pesos e ativações são quantizados de forma      
simulada (\textit{fake quantization}) durante o treinamento: na passagem \textit{forward}, os valores em ponto flutuante são projetados para os níveis discretos correspondentes       
($\hat{x}$), enquanto todos os cálculos e a etapa de \textit{backward} permanecem em ponto flutuante~\cite{nagel2021,pereira2023}. Dessa forma, os pesos são ajustados levando em 
consideração o erro de quantização inserido na função de perda, tornando o modelo mais resiliente e permitindo encontrar melhores soluções que o PTQ~\cite{nagel2021,pereira2023}.

Entretanto, durante o \textit{backpropagation}, é necessário calcular o gradiente da operação de arredondamento para o inteiro mais próximo ($\lfloor \cdot \rceil$), cuja derivada é  
nula quase em toda parte ou indefinida nos pontos de transição, o que inviabilizaria o treinamento baseado em gradiente~\cite{lin2023,nagel2021}. Para contornar essa limitação,
utiliza-se o \textit{Straight-Through Estimator} (STE), que aproxima a derivada do arredondamento por $1$ ($\partial \lfloor y \rceil / \partial y \approx 1$). Considerando a
saturação nos limiares $[q_{\min},\, q_{\max}]$, o gradiente do bloco quantizador em relação à sua entrada $x$ é dado por~\cite{nagel2021}:
\begin{equation} 
\frac{\partial \hat{x}}{\partial x} \approx
\begin{cases} 
1, & \text{se } q_{\min} \le x \le q_{\max}, \\ 
0, & \text{caso contrário.}
\end{cases}
\label{eq:ste} 
\end{equation}

Inicialmente, no QAT convencional, apenas os pesos e as ativações da rede eram ajustados por gradiente, enquanto os intervalos de quantização permaneciam fixos ou eram atualizados apenas por       
estatísticas, como \textit{Min-Max}. Contudo, as abordagens posteriores passaram a empregar o STE para otimizar também os parâmetros da quantização, $s$ e $\beta$, durante o           
treinamento~\cite{bhalgat2020,nagel2021}.

Para otimizar o fator de escala $s$ tanto nos pesos quanto nas ativações, surgiu o \textit{Learned Step Size Quantization} (LSQ), que adota a quantização simétrica ($\beta = 0$)      
definida em~\eqref{eq:quant_sim_q} e~\eqref{eq:quant_sim_dq}, porém tratando $s$ como um parâmetro treinável~\cite{bhalgat2020}. Dessa forma, durante o \textit{backpropagation}, o    
gradiente do quantizador em relação ao passo $s$ também é calculado via STE~\cite{bhalgat2020,nagel2021}:                                                                              
\begin{equation}                                                                                                                                                                       
\frac{\partial \hat{x}}{\partial s} \approx                                                                                                                                            
\begin{cases}                                                                                                                                                                          
-\dfrac{x}{s} + \left\lfloor \dfrac{x}{s} \right\rceil, & \text{se } q_{\min} \le x \le q_{\max}, \\                                                                                   
n, & \text{se } x < q_{\min}, \\                                                                                                                                                       
p, & \text{se } x > q_{\max}.                                                                                                                                                          
\end{cases}                                                                                                                                                                            
\label{eq:grad_lsq_s}                                                                                                                                                                  
\end{equation}                                                                                                                                                                         
Todavia, por adotar quantização simétrica, o LSQ perde eficiência quando os valores, especialmente ativações ou entradas, possuem distribuição assimétrica em torno de zero. Caso escolha um intervalo de quantização sem sinal, descartam-se os valores negativos, por outro lado escolher um intervalo com sinal divide os níveis igualmente entre os lados negativo e positivo,
desperdiçando resolução~\cite{bhalgat2020}. 

Para superar essa limitação, foi proposto o \textit{Learned Step Size Quantization Plus} (LSQ+), que adota a quantização assimétrica definida em~\eqref{eq:quant_assim_q}
e~\eqref{eq:quant_assim_dq}. Nesse método, além de treinar os pesos e o fator de escala $s$, ajusta-se também o parâmetro de deslocamento $\beta$ para lidar com distribuições         
assimétricas~\cite{bhalgat2020}. Com a inclusão de $\beta$, os gradientes em relação a $s$ e $\beta$ são obtidos via STE da seguinte forma~\cite{bhalgat2020}:
\begin{equation}
\frac{\partial \hat{x}}{\partial s} \approx
\begin{cases}
-\dfrac{x - \beta}{s} + \left\lfloor \dfrac{x - \beta}{s} \right\rceil, & \text{se } q_{\min} < x < q_{\max}, \\
n, & \text{se } x \le q_{\min}, \\
p, & \text{se } x > q_{\max},
\end{cases}
\label{eq:grad_lsqplus_s}
\end{equation}
\begin{equation}
\frac{\partial \hat{x}}{\partial \beta} \approx
\begin{cases}
0, & \text{se } q_{\min} < x < q_{\max}, \\
1, & \text{caso contrário.}
\end{cases}
\label{eq:grad_lsqplus_beta}
\end{equation}
Em ambos os casos, o STE é empregado para aproximar $\partial \bar{x}/\partial s$ e $\partial \bar{x}/\partial \beta$. Já para os pesos, mantém-se a quantização simétrica ($\beta =   
0$), pois observa-se que eles costumam se distribuir simetricamente em torno de zero. Com isso, a quantização assimétrica das ativações não gera custo computacional adicional na      
inferência, uma vez que o termo de deslocamento pode ser pré-calculado e incorporado ao viés (\textit{bias}) da camada~\cite{bhalgat2020}. Por fim, para reduzir a sensibilidade à     
inicialização e o impacto de \textit{outliers} em baixas precisões, os parâmetros iniciais $(s_{\text{init}},\, \beta_{\text{init}})$ são determinados pela minimização do             
MSE~\cite{bhalgat2020}.

\subsection{Formulação sob a Perspectiva de Otimização: LSQ+ Híbrido}

Ao utilizar a função de ativação \textit{Rectified Linear Unit} (ReLU), os valores de saída tornam-se positivos ou nulos, pois todas as entradas negativas são mapeadas 
para zero. Como no LSQ+ a inicialização por MSE e a atualização via STE operam de forma irrestrita sobre $s$ e $\beta$, o otimizador pode deslocar $\beta$ para fora da origem a fim de
reduzir o erro médio no restante do intervalo. Consequentemente, com $n = 0$, o limite inferior real $q_{\min} = \beta$ deixa de coincidir com zero~\cite{bhalgat2020}. Desse modo, o
valor real $0$ passa a ser representado com erro após a quantização, injetando um viés constante nos neurônios inativos~\cite{nagel2021} e criando um acoplamento mal-condicionado   
entre $s$ e $\beta$ durante o treinamento~\cite{antoniou2007,souza2026}.

Para evitar esse erro na ativação ReLU, pode-se formular o LSQ+ como um problema de otimização com restrições, garantindo que o parâmetro de deslocamento não viole a região viável,
isto é, impondo a restrição ativa $\beta = 0$~\cite{antoniou2007,souza2026}. Uma forma clássica de manter as iterações dentro do conjunto viável é o método de Projeção de Gradiente,  
que projeta a direção de busca sobre o hiperplano definido pelas restrições linearmente ativas~\cite{souza2026}. Nesse caso, como a restrição de fronteira exige $\beta = 0$ nas       
ativações pós-ReLU, a projeção mantém esse deslocamento nulo durante todo o treinamento enquanto os pesos e o fator de escala $s$ seguem sendo atualizados livremente.           
Ao mesmo tempo, mantém-se $\beta$ livre apenas onde existem valores negativos e positivos assimétricos, como na entrada normalizada do modelo, originando a variante LSQ+ Híbrido       
avaliada neste trabalho.

\section{Metodologia}

Conforme apresentado na introdução, este estudo toma como referência o trabalho ``Sistema Embarcado para Reconhecimento de Palavras usando ESP32-S3 e TinyML''~\cite{carlos2025}.      
Todavia, em vez da plataforma Edge Impulse, utilizou-se o PyTorch, mantendo rigorosamente o mesmo conjunto de dados e as mesmas configurações de pré-processamento, arquitetura e      
treinamento adotadas na plataforma original.

\subsection{Conjunto de Dados e Extração de Características (MFCC)}

O conjunto de dados totaliza 1.302~áudios, sendo 1.050 destinados ao desenvolvimento do modelo e 252 mantidos isolados para teste. Cada áudio foi rotulado em uma de seis classes      
(\textit{zero}, \textit{um}, \textit{dois}, \textit{três}, \textit{ruído} e \textit{desconhecido}) e padronizado temporalmente em 16~kHz com duração de 1~segundo (16.000~amostras),   
aplicando-se corte ou preenchimento com zeros (\textit{zero-padding}). Durante a etapa de desenvolvimento, as 1.050~amostras foram divididas em 80\% para treinamento (840~amostras) e 
20\% para validação (210~amostras)~\cite{carlos2025}.

Para extrair as características acústicas utilizadas no reconhecimento de voz, empregaram-se os Coeficientes Mel-Cepstrais (MFCC), técnica que modela a resposta não linear da audição 
humana --- mais sensível às baixas frequências do que às altas ---, seguida de normalização cepstral por média e variância (CMVN)~\cite{carlos2025,pereira2023}. Os parâmetros adotados
reproduzem integralmente a configuração do Edge Impulse e estão resumidos na Tabela~\ref{tab:mfcc}.

\begin{table}[htbp]
\centering
\caption{Parâmetros de extração MFCC e normalização CMVN.}
\label{tab:mfcc}
\begin{tabular}{lc} 
\toprule
\textbf{Parâmetro} & \textbf{Valor} \\ 
\midrule
Coeficientes cepstrais (MFCC) & 13 \\ 
Comprimento e passo do quadro & 20~ms (320 amostras) \\
Número de filtros Mel & 32 \\
Tamanho da FFT & 512~pontos \\
Frequência mínima & 10~Hz \\
Coeficiente de pré-ênfase & 0,98 \\
Janela de normalização CMVN & 101~quadros \\
\bottomrule
\end{tabular}
\end{table}

Ao final do pré-processamento, cada áudio gera uma matriz de entrada de dimensões $(13, 50)$, correspondendo a 13~coeficientes ao longo de 50~quadros temporais. No conjunto de
treinamento, os valores normalizados apresentaram média $-0{,}000$, desvio padrão $0{,}999$ e extremos assimétricos entre $-6{,}362$ e $+5{,}051$.

\subsection{Arquitetura da Rede Neural (1D-CNN)}

A rede neural convolucional unidimensional (1D-CNN) adotada neste trabalho reproduz a arquitetura empregada no Edge Impulse~\cite{carlos2025}, totalizando 1.974~parâmetros treináveis
distribuídos conforme detalhado na Tabela~\ref{tab:arquitetura}.                                                                                                                       

\begin{table}[htbp]
\centering
\setlength{\tabcolsep}{3.5pt}
\caption{Arquitetura da rede 1D-CNN (\textit{KWS\_CNN}) e número de parâmetros.}
\label{tab:arquitetura}
\begin{tabular}{lcc}   
\toprule            
\textbf{Camada} & \textbf{Saída} & \textbf{Parâmetros} \\
\midrule
Entrada MFCC & $(13, 50)$ & --- \\
\texttt{Conv1d} ($13 \to 8$, $k=3$, $p=1$) + ReLU & $(8, 50)$ & 320 \\
\texttt{MaxPool1d} ($k=2$, \textit{ceil}) + \texttt{Dropout} ($0{,}25$) & $(8, 25)$ & --- \\
\texttt{Conv1d} ($8 \to 16$, $k=3$, $p=1$) + ReLU & $(16, 25)$ & 400 \\
\texttt{MaxPool1d} ($k=2$, \textit{ceil}) + \texttt{Dropout} ($0{,}25$) & $(16, 13)$ & --- \\
\texttt{Flatten} + \texttt{Linear} ($208 \to 6$) & $(6)$ & 1.254 \\
\midrule
\textbf{Total} & --- & \textbf{1.974} \\
\bottomrule
\end{tabular}
\end{table}

\begin{figure}[htbp]
\centering
\resizebox{\columnwidth}{!}{\begin{tikzpicture}[
    font=\scriptsize,
    >=Stealth,
    node distance=0.24cm and 0.35cm,
    layer/.style={
        rectangle, rounded corners=2pt, draw=black!75, fill=gray!8,
        minimum width=3.9cm, minimum height=0.44cm, align=center, inner sep=2pt
    },
    quant/.style={
        rectangle, rounded corners=2pt, draw=blue!70!black, fill=blue!10,
        thick, minimum width=1.45cm, minimum height=0.40cm, align=center,
        font=\scriptsize\bfseries, inner sep=2pt
    },
    weight/.style={
        rectangle, rounded corners=2pt, draw=black!60, fill=white,
        minimum width=1.15cm, minimum height=0.40cm, align=center, inner sep=2pt
    },
    io/.style={
        rectangle, rounded corners=2pt, draw=black!80, fill=gray!18,
        minimum width=3.9cm, minimum height=0.44cm, align=center,
        font=\scriptsize\bfseries, inner sep=2pt
    }
]

% Fluxo principal (coluna esquerda)
\node[io]    (inp)   {Entrada MFCC $(13, 50)$};
\node[quant, below=of inp, minimum width=2.4cm] (qin) {Quantizador \texttt{in}};
\node[layer, below=of qin]  (conv1) {\texttt{Conv1d} ($13 \to 8$, $k=3$) + \texttt{ReLU}};
\node[quant, below=of conv1, minimum width=2.4cm] (qa1) {Quantizador \texttt{a1}};
\node[layer, below=of qa1]  (pool1) {\texttt{MaxPool1d} ($k=2$) + \texttt{Dropout} ($0{,}25$)};
\node[layer, below=of pool1] (conv2) {\texttt{Conv1d} ($8 \to 16$, $k=3$) + \texttt{ReLU}};
\node[quant, below=of conv2, minimum width=2.4cm] (qa2) {Quantizador \texttt{a2}};
\node[layer, below=of qa2]  (pool2) {\texttt{MaxPool1d} ($k=2$) + \texttt{Dropout} ($0{,}25$)};
\node[layer, below=of pool2] (flat)  {\texttt{Flatten} ($208$)};
\node[layer, below=of flat]  (fc)    {\texttt{Linear} ($208 \to 6$)};
\node[io,    below=of fc]    (out)   {Saída / \textit{Logits} (6 classes)};

% Ramos laterais de pesos
\node[quant, right=0.45cm of conv1] (qw1)  {Quant. \texttt{w1}};
\node[weight, right=0.30cm of qw1]  (w1)   {Pesos $w_1$};

\node[quant, right=0.45cm of conv2] (qw2)  {Quant. \texttt{w2}};
\node[weight, right=0.30cm of qw2]  (w2)   {Pesos $w_2$};

\node[quant, right=0.45cm of fc]    (qwfc) {Quant. \texttt{wfc}};
\node[weight, right=0.30cm of qwfc] (wfc)  {Pesos $w_{fc}$};

% Conexoes principais
\draw[->] (inp)   -- (qin);
\draw[->] (qin)   -- (conv1);
\draw[->] (conv1) -- (qa1);
\draw[->] (qa1)   -- (pool1);
\draw[->] (pool1) -- (conv2);
\draw[->] (conv2) -- (qa2);
\draw[->] (qa2)   -- (pool2);
\draw[->] (pool2) -- (flat);
\draw[->] (flat)  -- (fc);
\draw[->] (fc)    -- (out);

% Conexoes dos pesos
\draw[->] (w1)   -- (qw1);
\draw[->] (qw1)  -- (conv1);
\draw[->] (w2)   -- (qw2);
\draw[->] (qw2)  -- (conv2);
\draw[->] (wfc)  -- (qwfc);
\draw[->] (qwfc) -- (fc);

\end{tikzpicture}
}
\caption{Arquitetura \textit{KWS\_CNN} com os seis pontos de quantização simulada (\texttt{in}, \texttt{w1}, \texttt{a1}, \texttt{w2}, \texttt{a2} e \texttt{wfc}); os vieses (\textit{biases}) permanecem em FP32.}
\label{fig:arquitetura}
\end{figure}

Para avaliar as estratégias de quantização, a quantização simulada foi inserida em seis pontos da rede: na entrada MFCC (\texttt{in}), nos pesos das camadas convolucionais e da camada
densa (\texttt{w1}, \texttt{w2} e \texttt{wfc}) e nas ativações pós-ReLU (\texttt{a1} e \texttt{a2}), conforme ilustrado na Figura~\ref{fig:arquitetura}. A entrada \texttt{in} utiliza
representação com sinal e simétrica ($\beta = 0$) no PTQ, QAT Padrão e LSQ, passando para representação sem sinal com $\beta$ treinável no LSQ+ e no LSQ+ Híbrido. Os pesos    
(\texttt{w1}, \texttt{w2} e \texttt{wfc}) mantêm quantização com sinal e simétrica ($\beta = 0$) em todas as técnicas. Já as ativações pós-ReLU (\texttt{a1} e \texttt{a2}) adotam 
representação sem sinal em todos os métodos, mantendo $\beta = 0$ no PTQ, QAT Padrão, LSQ e LSQ+ Híbrido, e com $\beta$ treinável apenas no LSQ+ padrão. Os vieses (\textit{biases}) 
das camadas permanecem em ponto flutuante (FP32)~\cite{nagel2021}.

\subsection{Protocolo Experimental e Métricas de Avaliação}
% Treinamento baseline FP32 (300 epocas, Adam, lr=0,005), fine-tuning QAT (100 epocas, Adam, lr=1e-3, CosineAnnealingLR, seed 42), precisoes W8A8, W4A4 e W2A2, e metricas (Acuracia, Loss, ROC AUC e F1-Score ponderado).

\section{Resultados e Discussão}

% TODO: Apresentar as tabelas comparativas, graficos de convergencia, matrizes de confusao e analise critica.

\subsection{Desempenho Comparativo em 8, 4 e 2 Bits}
% Comparacao consolidada entre Baseline FP32, PTQ (Min-Max e MSE), QAT Padrao, QAT LSQ, QAT LSQ+ e QAT LSQ+ Hibrido em W8A8 (4x), W4A4 (8x) e W2A2 (16x).

\subsection{Análise de Otimização e Comportamento do LSQ+ Híbrido}
% Analise do desvio de ponto zero (zero-point drift) pos-ReLU no LSQ+ e efeito da projecao beta = 0 em a1/a2 mantendo beta treinavel na entrada MFCC.

\section{Conclusões}

% TODO: Sintetizar os principais achados em 8 bits (suficiencia do PTQ), 4 bits (ponto otimo do QAT LSQ) e 2 bits (colapso do PTQ Min-Max e recuperacao via QAT LSQ), bem como trabalhos futuros (deploy embarcado sub-8-bit no ESP32-S3, etc.).


\bibliographystyle{IEEEtran}
\bibliography{referencias}

\end{document}
